\section{Bidirectional Probe Lensing Hypothesis}
\label{sec:bidirectional-probe-lensing}

The preceding sections showed that the Dirac-, QFT-, and Einstein-type
regimes are not separate theories, but structural projections of one
common closure-stable master equation
(Proposition~\ref{prop:structural-lift-principle}). They also showed
that the Einstein reconstruction presently depends on a bilinear coframe
extraction whose nondegeneracy is still model-dependent. This suggests a
different route: rather than attempting to reconstruct geometry from a
single readout only, one may compare the traces of two probes approaching
the same structural core from different regime readings.

\subsection*{Basic setup}

Consider two probe modes of the triadic system,
$\Psi_{\mathrm{ART}}$ and $\Psi_{\mathrm{QFT}}$,
together with the orthogonal triadic component $\Psi_\perp$.
We regard $\Psi_{\mathrm{ART}}$ as a probe emphasizing the
large-scale/geometric regime and $\Psi_{\mathrm{QFT}}$ as a probe
emphasizing the phase-shifted observer/QFT regime at
$\zminus = +\bsym$ (Theorem~\ref{thm:symmetric-frame}).
The combined field is treated as a specialization of the triadic
two-probe system
\[
  \Psi_{\triangle}
  =
  \begin{pmatrix}
    \Psi_{\mathrm{ART}}\\
    \Psi_{\mathrm{QFT}}\\
    \Psi_\perp
  \end{pmatrix},
  \qquad
  \bigl(i\Gamma^\mu\mathcal{D}_\mu - \Omega_\triangle\bigr)
  \Psi_\triangle = 0.
\]
The third component is not decorative: it is the intrinsic coupling
channel through which the two regime-readings communicate.

\subsection*{Structural hypothesis}

\begin{proposition}[Bidirectional probe lensing hypothesis]
\label{hyp:bidirectional-probe-lensing}
\emph{(Hypothesis-grade; not yet theorem-level.)}
Let two probes approach the same closure-stable structural core from
different regime-readings, one on the ART side and one on the QFT side.
Then the compensator/closure core acts as an effective lensing kernel
for their traces. More precisely:
\begin{enumerate}[label=(\roman*)]
  \item the two probes need not follow identical local trajectories;
  \item their relative deflection, phase lag, and focusing structure are
        controlled by the same compensator/closure sector
        $C = \Pi^{(\mathcal{M})}_{\mathrm{spec}}$
        (Theorem~\ref{thm:hc-spectral-projector});
  \item the orthogonal triadic channel $\Psi_\perp$ mediates the
        nontrivial coupling between the two traces;
  \item the resulting trace geometry carries information about the
        effective metric readout before a full bilinear coframe
        reconstruction is imposed.
\end{enumerate}
This is a readout-relative statement
(Definition~\ref{def:readout-relative-lemma}): the lensing geometry is
visible in the bidirectional probe readout, while the kernel-level
content is the common closure-stable core.
\end{proposition}

\begin{theorem}[Theorem-level shadow of the bidirectional probe lensing hypothesis]
\label{thm:theorem-level-shadow-bidirectional-probe-lensing}
At the active theorem level, the present framework already supports the following weaker probe-lensing shadow:
\begin{enumerate}[leftmargin=2em,label=(\roman*)]
  \item two probe readouts approaching the same closure-stable structural core may be treated as distinct regime-readings of one common carrier rather than as probes of two disconnected foundations;
  \item the orthogonal triadic component is already a genuine coupling channel between the two regime-readings, not a decorative extra field;
  \item what is \emph{not} yet theorem-level is a full geometric law that identifies the resulting relative deflection, focusing, or phase-lag pattern with one explicit effective lensing tensor in complete generality.
\end{enumerate}
Therefore the proven theorem-level shadow of Hypothesis~\ref{hyp:bidirectional-probe-lensing} is:
\emph{the two probes already share one closure-stable carrier and one intrinsic coupling channel, even though a full lensing law for their relative geometry is not yet proved.}
\end{theorem}

\begin{proof}
The active framework already treats the ART- and QFT-side equations as distinguished regime/readout projections of one common closure-stable structural kernel rather than as disconnected sector theories. In the bidirectional probe setup, the combined triadic field
\[
  \Psi_{\triangle}
  =
  \begin{pmatrix}
    \Psi_{\mathrm{ART}}\\
    \Psi_{\mathrm{QFT}}\\
    \Psi_\perp
  \end{pmatrix}
\]
is explicitly organized so that the orthogonal component $\Psi_\perp$ is the intrinsic channel through which the two regime-readings communicate. Hence the theorem-level content already available is the existence of one common carrier and one nontrivial coupling channel between the two probes. What remains unproved is the stronger claim that their relative deflection/focusing/phase-lag pattern is already governed by one explicit universal lensing law at theorem level.
\end{proof}

\begin{corollary}[Operational reading of the probe-lensing hypothesis in the present v\docversion\ release line]
\label{cor:operational-reading-bidirectional-probe-lensing-v320}
Until a full effective lensing law is derived, later arguments may already use Theorem~\ref{thm:theorem-level-shadow-bidirectional-probe-lensing} as the active theorem-level substitute for the weakest part of Hypothesis~\ref{hyp:bidirectional-probe-lensing}: one may work with two probe readouts of one common closure-stable carrier coupled through the orthogonal triadic channel, while keeping the stronger geometric lensing law itself deferred. In the v3.20 crystal language, this means that two probe-facing domain realizations may already be treated as role-faces of one common carrier corridor, without claiming a complete universal focusing tensor.
\end{corollary}

\begin{proof}
This is the operational reformulation of Theorem~\ref{thm:theorem-level-shadow-bidirectional-probe-lensing}. The theorem already isolates the precise point at which common-carrier/two-readout coupling is available theorematically, while the stronger lensing law remains hypothesis-grade.
\end{proof}

\begin{remark}[Status of the compressed interface-to-optical bridge layer in the present release line]
\label{rem:status-of-the-first-bundled-theorem-level-shadow-package-v320}
\label{rem:status-of-the-first-v320-crystal-shadow-interface}
\label{rem:status-of-the-terminal-v320-working-interface-surface}
\label{rem:status-of-the-first-interface-level-defect-face-vanishing-step}
\label{rem:status-of-the-first-interface-level-fluid-readout-demotion-step}
\label{rem:status-of-the-first-interface-level-common-carrier-projection-step}
\label{rem:status-of-the-first-interface-level-probe-coupling-step}
\label{rem:status-of-the-first-bundled-interface-application-package}
\label{rem:status-of-the-first-bundled-interface-compatibility-step-in-v321}
\label{rem:status-of-the-first-op3-op5-bridge-test-on-the-terminal-interface}
\label{rem:status-of-the-first-no-second-carrier-step-in-v322}
\label{rem:status-of-the-first-partial-op3-op5-closure-theorem-in-v322}
\label{rem:status-of-the-first-restricted-lift-back-step-in-v322}
\label{rem:status-of-the-first-prism-readout-reduction-step-in-v323}
\label{rem:status-of-the-first-residual-spectral-diagnostic-step-in-v323}
\label{rem:status-of-the-first-parametric-spectral-fit-layer-in-v323}
\label{rem:status-of-the-first-spectral-calibration-step-in-v323}
\label{rem:status-of-the-first-overdetermined-spectral-consistency-step-in-v323}
\label{rem:status-of-the-first-effective-index-reconstruction-step-in-v323}
\label{rem:status-of-the-first-carrier-profile-correction-step-in-v323}
\label{rem:status-of-the-first-reduced-spectral-stability-step-in-v323}
\label{rem:status-of-the-first-bundled-optical-approximation-package-in-v323}
The present line no longer expands the long v3.20--v3.23 ladder step by step inside the active core. Instead, the theorem-level shadow package, the crystal--shadow/interface layer, the terminal interface application package, the first OP~3/OP~5 bridge/closure consequences, and the first optical approximation package are retained through one compressed interface-to-optical bridge layer. The detailed development remains recoverable from earlier candidates, while the active line keeps only the citation-equivalent surface needed for current proof use.
\end{remark}

\begin{definition}[Compressed interface-to-optical bridge token]
\label{def:v320-crystal-shadow-interface-object}
\label{def:reduced-prism-readout-sector}
\label{def:prism-readout-residual-profile}
\label{def:first-effective-carrier-dispersion-fit-ansatz}
\label{def:first-reconstructed-effective-refractive-profile}
\label{def:first-carrier-profile-correction-decomposition}
Let
\[
\mathfrak B^{\mathrm{int\text{-}opt}}
:=
\bigl(
\Sigma^{\mathrm{shadow}}_{\mathrm{v20}},
\mathfrak I_{\mathrm{term}},
\mathfrak R_{\mathrm{prism}},
\rho_{\mathrm{res}},
\mathcal F_{\mathrm{disp}},
 n^{\mathrm{rec}}_{\mathrm{eff}},
 \Delta n_{\mathrm{corr}}
\bigr)
\]
denote the compressed interface-to-optical bridge token, where the entries stand respectively for the bundled theorem-level shadow package, the terminal crystal--shadow interface object, the reduced prism-readout sector, the residual spectral profile, the first effective dispersion-fit ansatz, the reconstructed effective index profile, and its first carrier-profile correction decomposition. The active line uses this token as the compact carrier of all v3.20--v3.23 interface-to-optical data that are still needed in the present v\docversion\ release line.
\end{definition}

\begin{lemma}[Compressed spectral entry lemma on the interface-to-optical bridge]
\label{lem:first-prism-readout-lemma-on-the-retained-carrier-surface}
\label{lem:first-residual-spectral-diagnostic-lemma}
From the token $\mathfrak B^{\mathrm{int\text{-}opt}}$ one may recover, up to citation-equivalent proof use, the reduced prism sector together with the first residual spectral diagnostic entry data. In particular, any later argument that needs only the passage from terminal interface readout to residual spectral comparison may cite the present lemma instead of reopening the separate prism-readout and residual-profile lemmas.
\end{lemma}

\begin{proof}
The token $\mathfrak B^{\mathrm{int\text{-}opt}}$ retains exactly the interface-level theorematic shadow content, the terminal crystal/interface content, and the first optical approximation data that were previously introduced as separate stages. Since the active line uses these stages only through their retained terminal/citation-equivalent content, the reduced prism entry and the first residual spectral comparison layer may be read back from the single token without reopening the former intermediate exposition.
\end{proof}

\begin{proposition}[Compressed interface-to-optical bridge package]
\label{prop:first-bundled-theorem-level-shadow-package}
\label{prop:first-v320-crystal-shadow-compatibility-principle}
\label{prop:terminal-citation-rule-for-the-present-v320-working-interface}
\label{prop:first-interface-level-defect-face-vanishing-principle}
\label{prop:first-interface-level-fluid-readout-demotion-principle}
\label{prop:first-interface-level-common-carrier-projection-principle}
\label{prop:first-interface-level-probe-coupling-principle}
\label{prop:first-bundled-interface-application-package}
\label{prop:first-op3-op5-bridge-test-on-the-terminal-interface}
\label{prop:first-parametric-spectral-fit-principle}
\label{prop:first-three-parameter-spectral-calibration-principle}
\label{prop:first-overdetermined-spectral-consistency-principle}
\label{prop:first-effective-index-reconstruction-principle-from-prism-data}
\label{prop:first-carrier-profile-correction-principle}
\label{prop:first-reduced-spectral-stability-principle}
\label{prop:first-bundled-optical-approximation-package}
Assume a later argument depends only on the theorem-level shadow/common-carrier content, the terminal interface/crystal application layer, the weak OP~3/OP~5 bridge consequences already retained on that terminal layer, and the first optical approximation package given by reduced prism readout, residual spectral comparison, parametric fit, calibration, consistency, effective-index reconstruction, carrier-profile correction, and reduced spectral stability. Assume moreover that the argument does not require any still-deferred stronger claim such as an explicit HC meta-field equation, a deeper compensator generator, a second primary carrier, a full universal probe-lensing law, or a stronger optical inversion theorem beyond the first retained approximation package. Then the argument may work entirely with the single bridge token $\mathfrak B^{\mathrm{int\text{-}opt}}$ instead of reopening the former v3.20--v3.23 ladder step by step.
\end{proposition}

\begin{proof}
All formerly separate ingredients in the stated range were already used in the active line only through retained theorem-level or citation-level content: bundled theorematic shadows, terminal interface application, weak OP~3/OP~5 bridge consequences, and the first optical approximation package. The token $\mathfrak B^{\mathrm{int\text{-}opt}}$ was defined precisely to retain this active content and forget only the now-redundant intermediate exposition. Therefore any later argument that uses no stronger claim than the ones listed can cite the bridge token and stay on the compressed interface-to-optical surface.
\end{proof}

\begin{theorem}[Compressed theorematic consequences on the interface-to-optical bridge]
\label{thm:first-bundled-interface-compatibility-theorem}
\label{thm:first-no-second-carrier-theorem-on-the-terminal-interface}
\label{thm:first-partial-op3-op5-closure-theorem-on-the-terminal-interface}
\label{thm:first-restricted-lift-back-theorem-from-the-terminal-interface}
On the admissible theorematic window represented by $\mathfrak B^{\mathrm{int\text{-}opt}}$, the active line already retains, in compressed form, the compatibility theorem of the terminal interface package, the no-second-carrier consequence, the first partial OP~3/OP~5 closure consequence, and the first restricted lift-back consequence. Hence these theorematic uses may be cited through the single compressed theorematic bridge whenever no finer reopening below the retained interface-to-optical surface is needed.
\end{theorem}

\begin{proof}
The separate theorems in the former expanded presentation all concerned theorematic consequences that had already stabilized on the same terminal interface/application layer and that were then carried forward only through their retained citation content. Since the token $\mathfrak B^{\mathrm{int\text{-}opt}}$ preserves precisely that retained layer, the corresponding theorematic consequences may likewise be cited in compressed form on the same admissible window.
\end{proof}

\begin{corollary}[Compressed use rule for the interface-to-optical bridge]
\label{cor:bundled-use-rule-v320-theorem-level-shadow-package}
\label{cor:use-rule-for-the-v320-crystal-shadow-interface}
\label{cor:operational-mini-checklist-for-the-terminal-v320-working-interface}
\label{cor:interface-level-use-rule-for-exact-sector-defect-face-vanishing}
\label{cor:interface-level-use-rule-for-coarse-grained-fluid-readout}
\label{cor:interface-level-use-rule-for-common-carrier-and-projection-architecture}
\label{cor:interface-level-use-rule-for-probe-coupled-readout-architecture}
\label{cor:bundled-use-rule-for-the-terminal-interface-application-package}
\label{cor:bundled-use-rule-for-the-first-interface-compatibility-theorem}
\label{cor:use-rule-for-the-first-op3-op5-bridge-test}
\label{cor:use-rule-for-the-first-no-second-carrier-theorem}
\label{cor:use-rule-for-the-first-partial-op3-op5-closure-theorem}
\label{cor:use-rule-for-the-first-restricted-lift-back-theorem}
\label{cor:first-spectral-approximation-entry-rule-from-prism-readout}
\label{cor:first-spectral-fit-entry-rule-from-prism-residuals}
\label{cor:first-carrier-fit-workflow-for-optical-data}
\label{cor:first-identifiability-rule-for-optical-fit-parameters}
\label{cor:first-validation-workflow-for-measured-optical-series}
\label{cor:first-inversion-workflow-from-prism-angles-to-carrier-profile}
\label{cor:first-refinement-workflow-for-the-effective-carrier-law}
\label{cor:first-stopping-rule-for-the-reduced-prism-readout-fit-loop}
\label{cor:bundled-use-rule-for-the-first-optical-approximation-package}
Whenever a later argument needs only the active content compressed into $\mathfrak B^{\mathrm{int\text{-}opt}}$, it is sufficient to cite the present bridge layer and proceed directly to the present optical-diagnostic control surface. A finer reopening is required only when the argument depends on information intentionally forgotten by the compression: stronger generator-level structure, a deeper carrier split, a stronger universal lensing law, or a more detailed optical inversion/construction below the retained first approximation package.
\end{corollary}

\begin{proof}
This is the operational reformulation of Proposition~\ref{prop:first-bundled-optical-approximation-package} and Theorem~\ref{thm:first-partial-op3-op5-closure-theorem-on-the-terminal-interface} in their compressed bridge form. The present layer isolates the exact point at which the active line may stay on the retained interface-to-optical surface and pass directly into the present optical-diagnostic control machinery.
\end{proof}

\begin{remark}[Status of the compressed optical-diagnostic control layer in the present release line]
\label{rem:status-of-the-compressed-optical-diagnostic-control-layer-in-v324}
The present v\docversion\ release line now carries the whole first optical-to-control pipeline in compressed form: reference comparison, shortlist formation, structural back-reading, reopening priority, bombe/sieve reduction, tensor/shell control, early-depth policy, local utility ranking, and the first local spectral/QAM readout are retained as one finite citation layer rather than a long sequence of isolated micro-steps.
\end{remark}

\begin{definition}[Reference-comparison optical signature]
\label{def:reference-comparison-optical-signature}
For a reduced retained-carrier model series $S_{\mathrm{mod}}(\lambda_i)$ and a same-grid reference series $S_{\mathrm{ref}}(\lambda_i)$, define the comparison signature
\[
\Sigma_{\mathrm{opt}}(\lambda_i):=S_{\mathrm{ref}}(\lambda_i)-S_{\mathrm{mod}}(\lambda_i).
\]
\end{definition}

\begin{proposition}[First data-anchored optical comparison principle]
\label{prop:first-data-anchored-optical-comparison-principle}
Assume the retained-carrier model has already passed the bundled first-order optical validation and that $S_{\mathrm{ref}}$ is used only as an external same-observable comparison anchor. Then $\Sigma_{\mathrm{opt}}$ already supports the first trichotomy: uniformly small discrepancy gives first-order compatibility with the chosen reference family, smooth systematic discrepancy gives a shifted/refined placement inside that family, and irregular or localized discrepancy must be passed back to structural diagnostics before any stronger identification claim is made.
\end{proposition}

\begin{proof}
This is the compressed operational reading of the validated retained-carrier optical package: once like is compared with like and the reference is not promoted to a second carrier, the discrepancy can only be read as compatibility, smooth family shift, or a diagnostic reopen signal.
\end{proof}

\begin{definition}[Reference-family comparison score]
\label{def:reference-family-comparison-score}
For a finite family $\{S_{\mathrm{ref}}^{(\alpha)}\}_{\alpha\in\mathcal A}$ on the same wavelength grid, let
\[
\mathfrak s(\alpha):=\mathfrak s\bigl(\Sigma^{(\alpha)}_{\mathrm{opt}}\bigr)
\]
be any first-order score that decreases when the discrepancy is small, smooth, and structurally stable.
\end{definition}

\begin{proposition}[First reference-family ranking principle]
\label{prop:first-reference-family-ranking-principle}
Under the same retained-carrier assumptions, the finite family $\mathcal A$ may already be ranked by $\mathfrak s(\alpha)$, with lower score interpreted as better first-order compatibility. The output is not a final identification theorem but an ordered shortlist of admissible candidate families.
\end{proposition}

\begin{proof}
The comparison principle yields a controlled discrepancy profile for each candidate family; a finite score extraction therefore produces an admissible first shortlist without changing the carrier-side framework.
\end{proof}

\begin{corollary}[First shortlist workflow for candidate optical families]
\label{cor:first-shortlist-workflow-for-candidate-optical-families}
The first optical workflow is: validate the reduced model series, compare it with each reference family, score the resulting discrepancies, and retain the lowest-scoring families as the first optical shortlist.
\end{corollary}

\begin{proof}
This is the direct operational form of Propositions~\ref{prop:first-data-anchored-optical-comparison-principle} and \ref{prop:first-reference-family-ranking-principle}.
\end{proof}

\begin{definition}[Optical shortlist back-reading pattern]
\label{def:optical-shortlist-back-reading-pattern}
A back-reading pattern is the first structural interpretation of the shortlist class, distinguishing smooth boundary-dominated, smooth carrier-dominated, smooth mixed boundary/carrier, and localized defect/readout-near mismatch patterns.
\end{definition}

\begin{proposition}[First structural back-reading principle from optical shortlists]
\label{prop:first-structural-back-reading-principle-from-optical-shortlists}
Every finite shortlist obtained from Corollary~\ref{cor:first-shortlist-workflow-for-candidate-optical-families} may already be read back into one of the first structural classes of Definition~\ref{def:optical-shortlist-back-reading-pattern}. Smooth uniform offset points to boundary reopening, smooth curvature drift points to carrier reopening, coexistence of both points to mixed reopening, and localized/sign-changing mismatch points to defect/readout reopening.
\end{proposition}

\begin{proof}
This is exactly the first inverse reading of the retained optical discrepancy: the geometric shape of the residual determines which structural sector must be reopened first.
\end{proof}

\begin{corollary}[First workflow from optical shortlist to structural diagnostic hint]
\label{cor:first-workflow-from-optical-shortlist-to-structural-diagnostic-hint}
The first shortlist-to-structure workflow is: build the shortlist, classify its residual shape, and pass the result to the corresponding structural reopening hint.
\end{corollary}

\begin{proof}
Immediate from Proposition~\ref{prop:first-structural-back-reading-principle-from-optical-shortlists}.
\end{proof}

\begin{definition}[Structural diagnostic reopening priority]
\label{def:structural-diagnostic-reopening-priority}
The reopening priority is the first finite ordering of admissible structural sectors induced by the back-reading class.
\end{definition}

\begin{proposition}[First diagnostic reopening-priority principle from optical back-reading]
\label{prop:first-diagnostic-reopening-priority-principle-from-optical-back-reading}
The first back-reading class determines the first reopening order: boundary first for smooth offset, carrier first for smooth curvature, mixed boundary/carrier first for smooth coexistence, and defect/readout first for localized or parity-sensitive triggers.
\end{proposition}

\begin{proof}
The back-reading class already singles out which sector best matches the observed discrepancy, so the coarse reopening order follows immediately.
\end{proof}

\begin{corollary}[First refinement workflow after optical back-reading]
\label{cor:first-refinement-workflow-after-optical-back-reading}
After optical back-reading, it is enough to open the leading sector first and postpone deeper alternative sectors unless the leading residual test fails.
\end{corollary}

\begin{proof}
Immediate from Definition~\ref{def:structural-diagnostic-reopening-priority} and Proposition~\ref{prop:first-diagnostic-reopening-priority-principle-from-optical-back-reading}.
\end{proof}

\begin{definition}[Structural bombe on a reopening menu]
\label{def:structural-bombe-on-a-reopening-menu}
Given a finite reopening menu $\mathcal M$, the structural bombe is the first elimination operator that removes candidates incompatible with the observed back-reading class, residual improvement test, and low-order admissibility constraints.
\end{definition}

\begin{proposition}[First structural bombe elimination principle on reopening menus]
\label{prop:first-structural-bombe-elimination-principle-on-reopening-menus}
On any finite reopening menu, the bombe eliminates every candidate that contradicts the active back-reading class or worsens the residual while preserving any candidate that stays structurally admissible and non-worsening.
\end{proposition}

\begin{proof}
This is the compressed elimination rule built into the first bombe layer: the menu shrinks by contradiction and non-improvement.
\end{proof}

\begin{definition}[Nested structural bombe]
\label{def:nested-structural-bombe}
A nested bombe is the same elimination mechanism applied to a finite submenu below a surviving parent branch.
\end{definition}

\begin{proposition}[Nested bombe backtracking principle]
\label{prop:nested-bombe-backtracking-principle}
If every child of a surviving parent branch is eliminated, then the parent branch itself must be reopened or discarded; otherwise the surviving child becomes the ordered next refinement path.
\end{proposition}

\begin{proof}
This is the finite backtracking rule for nested menus.
\end{proof}

\begin{definition}[First smooth mixed shortlist case]
\label{def:first-smooth-mixed-shortlist-case}
The first smooth mixed case is the shortlist regime where smooth offset and smooth curvature coexist, no leading localized defect trigger is visible, and the coarse menu contains the candidates $\mu_{\mathrm B},\mu_{\mathrm C},\mu_{\mathrm{BC}},\mu_{\mathrm D}$.
\end{definition}

\begin{proposition}[First bombe reduction in the smooth mixed optical case]
\label{prop:first-bombe-reduction-in-the-smooth-mixed-optical-case}
In the first smooth mixed case, the bombe eliminates $\mu_{\mathrm D}$ by absence of a defect trigger and eliminates the pure boundary and pure carrier branches whenever each leaves the complementary smooth residual visible. The surviving coarse branch is then $\mu_{\mathrm{BC}}$, and the first surviving child order is boundary correction before carrier-curvature refinement.
\end{proposition}

\begin{proof}
This is the first applied mixed-case elimination: smooth mixed evidence kills the defect branch and the two pure branches, leaving the ordered mixed refinement path.
\end{proof}

\begin{definition}[First defect/readout-near shortlist case]
\label{def:first-defect-readout-near-shortlist-case}
The first defect/readout-near case is the shortlist regime where the mismatch is localized or parity-sensitive and cannot be absorbed by smooth offset or smooth curvature alone.
\end{definition}

\begin{definition}[Readout-trigger sieve]
\label{def:readout-trigger-sieve}
The readout-trigger sieve is the first pre-bombe filter that retains only those candidates whose local trigger profile matches the observed localized defect/readout pattern.
\end{definition}

\begin{proposition}[First trigger-sieve-plus-bombe reduction in the defect/readout-near optical case]
\label{prop:first-trigger-sieve-plus-bombe-reduction-in-the-defect-readout-near-optical-case}
In the first defect/readout-near case, the trigger sieve removes the smooth candidates before the bombe stage, and the reduced bombe then retains the direct defect/readout branch whenever hybrid smooth-plus-readout branches fail to improve the localized trigger beyond the direct patch.
\end{proposition}

\begin{proof}
Localized trigger information is stronger than a smooth menu reading, so the sieve-first regime compresses the menu before the bombe makes the final coarse elimination.
\end{proof}

\begin{definition}[Minimal structural diagnostic transition tensor]
\label{def:minimal-structural-diagnostic-transition-tensor}
The minimal transition tensor is the four-axis array
\[
\Theta_{r,u,c}^{\;r'},
\]
with regime $r$, tool $u$, readout class $c$, and next regime $r'$, supported only on admissible structural moves.
\end{definition}

\begin{definition}[First diagnostic shell condition for the minimal transition tensor]
\label{def:first-diagnostic-shell-condition-for-the-minimal-transition-tensor}
A transition lies on the diagnostic shell if and only if it is simultaneously closure-compatible, transport-compatible, readout-compatible, and depth-compatible. Off-shell entries of $\Theta$ vanish.
\end{definition}

\begin{definition}[First explicit mini-assignment for the leading optical regimes]
\label{def:first-explicit-mini-assignment-for-the-leading-optical-regimes}
The first explicit mini-assignment sets the leading on-shell tool choice to bombe-first in the smooth mixed regime and sieve-first in the defect/readout-near regime.
\end{definition}

\begin{proposition}[Priority readout of the first explicit mini-assignment]
\label{prop:priority-readout-of-the-first-explicit-mini-assignment}
The mini-assignment already fixes the first local tool priority for the two leading regimes and therefore compresses the first branching node to a singleton tool choice.
\end{proposition}

\begin{proof}
This is the direct reading of the nonzero shell-supported entries of the mini-assignment.
\end{proof}

\begin{corollary}[What the explicit mini-assignment already buys for proof economy]
\label{cor:what-the-explicit-mini-assignment-already-buys-for-proof-economy}
Once the regime and readout class are identified, the first tool step may be cited directly from the mini-assignment instead of reopening the whole first tool menu.
\end{corollary}

\begin{proof}
Immediate from Proposition~\ref{prop:priority-readout-of-the-first-explicit-mini-assignment}.
\end{proof}

\begin{definition}[First explicit nested submenu assignment after the two leading entries]
\label{def:first-explicit-nested-submenu-assignment-after-the-two-leading-entries}
The first nested submenu assignment fixes the next ordered submenu move: boundary-offset repair before carrier-curvature repair in the smooth mixed line, and local defect/readout repair before any smooth supplement in the defect/readout-near line.
\end{definition}

\begin{proposition}[Second-step priority readout of the first nested submenu assignment]
\label{prop:second-step-priority-readout-of-the-first-nested-submenu-assignment}
The nested submenu assignment already compresses the second branching node to a singleton child priority on the two leading lines.
\end{proposition}

\begin{proof}
This is the second-step reading of the surviving nested submenu entries.
\end{proof}

\begin{definition}[First shell-supported branch-depth policy]
\label{def:first-shell-supported-branch-depth-policy}
The early branch-depth policy is the shell-supported rule that the first two branching depths are determined by the explicit menu assignment and the explicit nested submenu assignment.
\end{definition}

\begin{proposition}[Early-depth support propagation under the first branch-depth policy]
\label{prop:early-depth-support-propagation-under-the-first-branch-depth-policy}
Within the stabilized early window, support propagates exactly along the policy-selected branch and no competing sibling remains admissible.
\end{proposition}

\begin{proof}
This is the finite propagation rule induced by the first two singleton assignments.
\end{proof}

\begin{definition}[First shell-compatible local utility functional]
\label{def:first-shell-compatible-local-utility-functional}
Let $U$ be a local utility functional on on-shell candidates, combining residual gain, closure-compatibility, transport-compatibility, readout fidelity, and backtracking cost. Off-shell candidates are assigned utility $-\infty$.
\end{definition}

\begin{proposition}[Utility refinement of the early branch-depth policy]
\label{prop:utility-refinement-of-the-early-branch-depth-policy}
On the stabilized early window, the unique utility maximizer reproduces the existing singleton policy; beyond that window, utility provides the first ranked continuation rule on the remaining on-shell menu.
\end{proposition}

\begin{proof}
Within the early window the shell and the previous singleton policy already select one surviving move; utility therefore agrees there and extends the same ordering principle to later local menus.
\end{proof}

\begin{corollary}[What the first local utility layer already buys for later depth extension]
\label{cor:what-the-first-local-utility-layer-already-buys-for-later-depth-extension}
Later local continuation may be treated as shell-first, utility-second, and only afterwards as a hard singleton policy when stability persists.
\end{corollary}

\begin{proof}
Immediate from Proposition~\ref{prop:utility-refinement-of-the-early-branch-depth-policy}.
\end{proof}

\begin{definition}[First utility-guided provisional depth-$2$ continuation rule]
\label{def:first-utility-guided-provisional-depth-2-continuation-rule}
The provisional depth-$2$ rule ranks the admissible children of the first stabilized branch by the local utility functional.
\end{definition}

\begin{proposition}[Utility-supported continuation beyond the stabilized early-depth window]
\label{prop:utility-supported-continuation-beyond-the-stabilized-early-depth-window}
Beyond the first two stabilized depths, the next local move may be taken from the highest-ranked on-shell child so long as no break trigger, shell loss, or utility tie forces a reopening.
\end{proposition}

\begin{proof}
This is the first later-depth continuation principle induced by the utility layer.
\end{proof}

\begin{definition}[First explicit smooth-mixed depth-$2$ test case]
\label{def:first-explicit-smooth-mixed-depth-2-test-case}
In the smooth mixed line, the first depth-$2$ test case is the surviving mixed path after boundary-offset repair, with carrier-side curvature correction as the leading remaining local move.
\end{definition}

\begin{proposition}[First explicit third-step priority in the smooth mixed line]
\label{prop:first-explicit-third-step-priority-in-the-smooth-mixed-line}
The explicit smooth mixed depth-$2$ test case selects carrier-curvature follow-up as the unique third step.
\end{proposition}

\begin{proof}
This is the leading continuation left after the earlier mixed-case eliminations and utility ranking.
\end{proof}

\begin{definition}[First explicit defect/readout-near depth-$2$ test case]
\label{def:first-explicit-defectreadout-near-depth-2-test-case}
In the defect/readout-near line, the first depth-$2$ test case is the surviving direct readout path after the first local defect patch, with continued local defect/readout stabilization as the leading next move.
\end{definition}

\begin{proposition}[First explicit third-step priority in the defect/readout-near line]
\label{prop:first-explicit-third-step-priority-in-the-defectreadout-near-line}
The explicit defect/readout-near depth-$2$ test case selects continued local defect/readout stabilization as the unique third step.
\end{proposition}

\begin{proof}
This is the direct continuation of the sieve-first/readout-first line once the localized trigger remains the dominant residual feature.
\end{proof}

\begin{definition}[First QAM-coded local spectral menu operator]
\label{def:first-qam-coded-local-spectral-menu-operator}
For a finite shell-admissible local menu, define the diagonal QAM-coded menu operator whose diagonal entries are the corresponding local utility values.
\end{definition}

\begin{proposition}[Spectral singleton readout of shell-supported local continuation]
\label{prop:spectral-singleton-readout-of-shell-supported-local-continuation}
The finite spectral theorem for the local QAM-coded menu operator reproduces the utility winner as the top eigendirection whenever the local utility maximizer is unique.
\end{proposition}

\begin{proof}
A finite diagonal self-adjoint menu operator has the utility values as spectrum, so the unique maximizer is exactly the top local spectral readout.
\end{proof}

\begin{corollary}[What the first local spectral/QAM bridge already buys]
\label{cor:what-the-first-local-spectralqam-bridge-already-buys}
Local continuation can now be cited either as the unique utility winner or as the unique top eigendirection of the corresponding local QAM-coded menu operator.
\end{corollary}

\begin{proof}
Immediate from Proposition~\ref{prop:spectral-singleton-readout-of-shell-supported-local-continuation}.
\end{proof}
\begin{remark}[Status of the compressed later-depth controller surface in the present release line]
\label{rem:status-of-the-compressed-laterdepth-controller-surface-in-v324}
The detailed later-depth controller build-up of the preceding line can now be treated in compressed form on the active line. Instead of carrying the full micro-chain of propagation, break, reentry, decoder, action, consistency, patch, decomposition, concatenation/cut, and finite normal form as separate local interfaces, the present line packages the same finite admissible later-depth content into one exact citation surface.
\end{remark}

\begin{definition}[Compressed later-depth controller surface on admissible windows]
\label{def:compressed-laterdepth-controller-surface-on-admissible-windows}
Let
\[
I=[d_-,d_+]\cap\mathbb N
\]
be a finite admissible later-depth interval on the leading optical regimes. We say that $I$ is read on the \emph{compressed later-depth controller surface} if its later-depth control content is cited only through the token
\[
\mathfrak Z^{\mathrm{cyc}}(I)
=
\langle \mathsf{laterdepth},d_-,d_+,\mathfrak N^{\mathrm{cyc}}(I)\rangle,
\]
where $\mathfrak N^{\mathrm{cyc}}(I)$ is the finite controller normal form storing the ordered lengths of maximal propagated advance-patches together with the intervening separator symbols $\mathsf{reopen}$ and $\mathsf{anchor}$.

Equivalently, the compressed surface remembers exactly the finite alternation pattern
\[
\text{advance-patch}\; /\; \text{separator}\; /\; \text{advance-patch}\; /\; \cdots
\]
on $I$ and forgets only the inessential internal relabeling inside propagated patch interiors.
\end{definition}

\begin{proposition}[Exact finite control recovery from the compressed later-depth controller surface]
\label{prop:exact-finite-control-recovery-from-the-compressed-laterdepth-controller-surface}
Assume Definition~\ref{def:compressed-laterdepth-controller-surface-on-admissible-windows}. Let $I$ be a finite admissible later-depth interval.
\begin{enumerate}[leftmargin=2em,label=\textup{(\roman*)}]
  \item The token $\mathfrak Z^{\mathrm{cyc}}(I)$ determines the exact break-aware later-depth controller profile of $I$ up to internal depth-preserving relabeling inside maximal propagated advance-patches.
  \item Every later argument whose dependence on $I$ is only through that finite control profile may be written entirely on the compressed later-depth controller surface.
  \item Local gluing and cutting of adjacent admissible windows are recoverable on the compressed surface by concatenating and cutting the stored normal-form word at the corresponding boundary depth.
  \item Whenever a finite later-depth passage is used only through its controller pattern, citing $\mathfrak Z^{\mathrm{cyc}}(I)$ is exact and reopening the full micro-chain is unnecessary.
\end{enumerate}
\end{proposition}

\begin{proof}
The finite controller normal form already stores the unique alternating decomposition of $I$ into maximal propagated advance-patches and separator sites. The compressed token records exactly that normal form together with the interval endpoints, hence it carries the same finite controller content as the detailed later-depth micro-chain modulo the inessential relabeling inside propagated patch interiors. Therefore every argument depending only on this finite control pattern may be written against $\mathfrak Z^{\mathrm{cyc}}(I)$ alone. The local concatenation/cut behavior is recovered because gluing and cutting act directly on the same ordered normal-form word.
\end{proof}

\begin{corollary}[What the compressed later-depth controller surface already buys on the active line]
\label{cor:what-the-compressed-laterdepth-controller-surface-already-buys-on-the-active-line}
A substantial part of the later-depth controller apparatus is now available in one compact active citation layer. In particular, later arguments may cite a finite admissible later-depth window through
\[
\mathfrak Z^{\mathrm{cyc}}(I)
\]
instead of reopening the full sequence of propagation, break, reentry, decoder, action, consistency, controller, patch, decomposition, and concatenation/cut statements.
\end{corollary}

\begin{proof}
This is the immediate operational reading of Proposition~\ref{prop:exact-finite-control-recovery-from-the-compressed-laterdepth-controller-surface}.
\end{proof}

\begin{remark}[Status of the terminal interface-application citation token in the present v\docversion\ release line]
\label{rem:status-of-the-terminal-interface-application-surface-in-the-present-v324-line}
The present line still does \emph{not} claim one unrestricted final application theory covering every later use of the framework. It does, however, support one compressed terminal citation layer: whenever a later argument uses only the currently bundled interface-application package on one terminal crystal--shadow interface, that package may be cited through one terminal token instead of reopening the individual interface-level applications.
\end{remark}

\begin{definition}[Compressed terminal interface-application token]
\label{prop:terminal-citation-rule-for-the-present-v324-interface-application-line}
Let
\[
\mathfrak I_{\mathrm{v20}}
\bigl([\mathfrak C_1,\dots,\mathfrak C_n]_{\star_{\mathrm{cr}}\mathrm{CNF}}\bigr)
\]
be a terminal interface object on the present release line. We write
\[
\mathfrak T^{\mathrm{app}}(\mathfrak I)
=
\langle
\mathsf{terminal},
\mathfrak I,
\mathsf{defectface},
\mathsf{fluiddem},
\mathsf{carrierproj},
\mathsf{twoprobe}
\rangle
\]
if the later argument depends only on the bundled interface-application package of Proposition~\ref{prop:first-bundled-interface-application-package}, requires no still-deferred stronger claim, and uses no pre-interface data already forgotten by passage to the terminal surface.
\end{definition}

\begin{proposition}[Exact terminal citation on the compressed interface-application token]
\label{cor:operational-mini-checklist-for-the-terminal-interface-application-surface}
Assume Definition~\ref{prop:terminal-citation-rule-for-the-present-v324-interface-application-line}. If a later argument depends on a terminal interface object only through the four bundled interface-level applications encoded in \(\mathfrak T^{\mathrm{app}}(\mathfrak I)\), then citing \(\mathfrak T^{\mathrm{app}}(\mathfrak I)\) is exact and reopening the individual interface-level application principles is unnecessary. If any stronger generator, deeper compensator, explicit fluid coarse-graining claim, universal lensing claim, or forgotten pre-interface datum is needed, the proof must reopen the earliest stage where that information is still visible.
\end{proposition}

\begin{proof}
This is the bundled interface-application package of Proposition~\ref{prop:first-bundled-interface-application-package} rewritten as one terminal citation token. Exactness holds because the token stores precisely the four currently stabilized interface-level uses and nothing stronger; if a later argument needs more, it is by definition outside the token and must reopen the earlier stage carrying the missing information.
\end{proof}











\begin{remark}
This hypothesis does not identify the compensator literally with a
classical gravitational lens. Its weaker and structurally safer claim
is that the closure/HC core behaves as an effective focusing kernel for
probe traces --- the spectral-projector image $\mathrm{im}(C) =
\mathcal{S}_0$ of Theorem~\ref{thm:hc-spectral-projector} acts as the
common structural center around which both probe traces are organized.
\end{remark}

\subsection*{Why this is plausible in the present framework}

Three already established facts support the hypothesis.

First, the master equation admits distinguished structural projections to
the probe/Dirac, phase-shifted QFT, and sharpened Einstein regimes
(Proposition~\ref{prop:structural-lift-principle}). Hence the ART and QFT
probes are not foreign objects but two regime readings of one common
structural dynamics.

Second, the triadic two-probe system already contains an intrinsic third
channel $\Psi_\perp$ coupling the two probe modes. The middle sector
therefore carries genuine interaction data, not a mere bookkeeping slot.

Third, the Einstein reconstruction is already expressed through bilinears
of a closed triadic probe field, so probe traces are not external
diagnostics but the natural source of geometric data in the current
framework.

\subsection*{Trace observables}

\paragraph*{Relative deflection.}
Let $X^\mu_{\mathrm{ART}}(\tau)$ and $X^\mu_{\mathrm{QFT}}(\tau)$
denote effective probe traces. The deflection observable is
\[
  \Delta X^\mu(\tau)
  :=
  X^\mu_{\mathrm{ART}}(\tau) - X^\mu_{\mathrm{QFT}}(\tau).
\]

\paragraph*{Relative phase lag.}
\[
  \Delta\phi(\tau)
  :=
  \phi_{\mathrm{ART}}(\tau) - \phi_{\mathrm{QFT}}(\tau).
\]
This quantity is sensitive to compensator-mediated observer-frame
displacement (Remark~\ref{rem:frame-shift-spectral}).

\paragraph*{Cross-bilinear geometry channel.}
\[
  E^{a}_{\mu,\mathrm{cross}}
  :=
  \operatorname{Re}\bigl(
    \bar\Psi_{\mathrm{ART}}
    \,\Gamma^a\,
    \mathcal{D}_\mu
    \Psi_{\mathrm{QFT}}
  \bigr),
\]
and similarly with probes exchanged. These mixed quantities detect
geometry encoded not in a single probe sector but in the coupled
relation of both probes through the triadic core.

\paragraph*{Focusing functional.}
\[
  \mathcal{L}_{\mathrm{focus}}
  :=
  f\!\bigl(
    E^{a}_{\mu,\mathrm{cross}},\,
    \mathfrak{R}_C,\,
    \mathcal{H}_C
  \bigr),
\]
where $f$ detects concentration of the two traces toward the same
closure-stable branch. The expression is heuristic at this stage.

\subsection*{Expected regimes}

\begin{proposition}[Expected qualitative behaviour of probe lensing]
\label{prop:expected-probe-lensing-behaviour}
The bidirectional probe lensing hypothesis predicts:
\begin{enumerate}[label=(\roman*)]
  \item under exact closure, the trace structure sharpens and becomes
        sector-clean;
  \item when the triadic coupling channel is switched off, the
        cross-bilinear focusing signal weakens or disappears;
  \item in the phase-shifted observer regime, relative phase lag is more
        pronounced than geometric coframe data;
  \item in the large-scale closure-stable regime, relative focusing is
        expected to organize into an effective geometric readout.
\end{enumerate}
\end{proposition}

\begin{proof}
Exact closure suppresses leakage and makes the physical sector
autonomous (Lemma~\ref{lem:closure-forcing-lemma}), so trace observables
become sector-clean. The triadic two-probe equation shows that
$\Psi_\perp$ is the intrinsic coupling route; without it the two regime
probes largely decouple, weakening the cross-bilinear signal. By
Proposition~\ref{prop:structural-lift-principle}, QFT emerges in the
phase-shifted observer regime while the Einstein limit appears only at
large closure-stable scales, so phase effects dominate the QFT reading
and geometric focusing emerges only later.
\end{proof}

\subsection*{Relation to Einstein reconstruction}

The bidirectional probe picture may weaken the hardest current
bottleneck of the Einstein block. At present, the metric is reconstructed
from the bilinear coframe
\[
  E^a{}_\mu
  =
  \operatorname{Re}\bigl(
    \bar\Psi_\triangle\,\Gamma^a\,\mathcal{D}_\mu\Psi_\triangle
  \bigr),
\]
and the unresolved step is to justify its maximal rank from the triadic
probe dynamics itself. The bidirectional probe hypothesis suggests an
earlier readout layer: geometry may first appear as a stable trace
focusing/deflection pattern before it is compressed into a single
effective coframe.

\begin{remark}
If this route works, the Einstein sector would no longer require a
direct leap from closure to metric. The structural lift
(Definition~\ref{def:structural-lift}) of this two-step chain is:
\[
\begin{aligned}
\text{closure-stable probe traces}
&\Longrightarrow \text{effective focusing geometry}\\
&\Longrightarrow \text{coframe / metric reconstruction}.
\end{aligned}
\]
This could reduce the present model-dependence of the Einstein limit and
would realize the concrete form of the meta-field projection
(iv) in Hypothesis~\ref{hyp:hc-meta-generator}.
\end{remark}

% ------------------------------------------------------------------
% Theorem 7: HC-Induced Readout Geometry (v2.27.4 integration)
% ------------------------------------------------------------------

